% Lecture example set: SC2008-L02
% Source: M1-L3 p. 17; M1-L4 p. 16
% Human verified: Yes

\exercisemeta
  {SC2008-L02-EX}
  {2026-08-18}
  {M1-L3 p.~17；M1-L4 p.~16}
  {两道 lecture examples 的答案已完成并核对}
  {已确认（2026-08-18）}

\exercisequestion{SC2008-L02-E01}{Stop-and-Wait Link Utilization}

\textbf{对应知识点：}
\relatedtopic{SC2008-L02-T02}{Stop-and-Wait 与 Sliding Window Flow Control}

\subsubsection*{题目}

一条 point-to-point link（点到点链路）的 transmission rate 为 $2.4$ Mbps，距离为 $50$ km，signal propagation velocity 为 $2\times10^8$ m/s，frame length 为 $300$ bytes。忽略 frame error、ACK transmission time 与 processing time，求 Stop-and-Wait link utilization（链路利用率）。

\begin{workedsolution}
先统一单位：$L=300\times8=2400$ bits。于是
\begin{align*}
  T_{\mathrm{frame}}
    &=\frac{2400}{2.4\times10^6}
      =1000\ \mu\mathrm{s},\\
  T_{\mathrm{prop}}
    &=\frac{50\times10^3}{2\times10^8}
      =250\ \mu\mathrm{s},\\
  a&=\frac{250}{1000}=0.25.
\end{align*}
因此
\[
  U=\frac{1}{1+2(0.25)}=\frac{2}{3}\approx66.7\%.
\]
若进一步计算实际 throughput，则
\[
  2.4\times\frac{2}{3}=1.6\ \mathrm{Mbps}.
\]
\end{workedsolution}

\textbf{检查点：}$T_{\mathrm{frame}}=L/R$ 描述把整个 frame 推入链路的时间；$T_{\mathrm{prop}}=d/v$ 描述信号跨越链路的时间，二者不能混用。

\exercisequestion{SC2008-L02-E02}{Stop-and-Wait ARQ with Frame Loss}

\textbf{对应知识点：}
\relatedtopic{SC2008-L02-T04}{Stop-and-Wait ARQ}

\subsubsection*{题目}

沿用 E01 的 link parameters（链路参数），并令一次 frame transmission 的 loss probability（丢失概率）为 $P=0.1$。求 Stop-and-Wait ARQ 的 link utilization。

\begin{workedsolution}
由 E01 得 $a=0.25$。因此
\[
  U_{\mathrm{Stop-and-Wait\ ARQ}}
  =\frac{1-P}{1+2a}
  =\frac{0.9}{1.5}
  =0.6.
\]
若进一步计算实际 throughput，则
\[
  2.4\times0.6=1.44\ \mathrm{Mbps}.
\]
\end{workedsolution}

\textbf{检查点：}$1-P$ 是每次独立 transmission attempt 成功的长期比例；若 retransmission 仍可能失败，就不能用 $1/(1+P)$ 替代。
